\Needspace{7\baselineskip}\Article{298}
\textsc{Problem.} \textit{Let \(a,b,c\) be any numbers, none of them \(=0\). Find the conditions under which the equation}

\[axx+byy+czz=0\ldots(\omega).\]
\textit{Solution.} Let \(\alpha\alpha,\beta\beta,\gamma\gamma\) be the largest squares dividing \(bc,ac,ab\), respectively, and set \(\alpha a=\beta\gamma A,\ \beta b=\alpha\gamma B,\ \gamma c=\alpha\beta C\). Then \(A,B,C\) will be pairwise coprime integers. The equation \((\omega)\) will be solvable or unsolvable according as
\[AXX+BYY+CZZ=0\ldots(\Omega)\]
is solvable or unsolvable. This can be decided by Article 294.

\textit{Proof.} Set \(bc=\mathfrak A\alpha\alpha,\ ac=\mathfrak B\beta\beta,\ ab=\mathfrak C\gamma\gamma\). Then \(\mathfrak A,\mathfrak B,\mathfrak C\) are integers free of square factors, and \(\mathfrak A=BC,\ \mathfrak B=AC,\ \mathfrak C=AB\). So \(\mathfrak A\mathfrak B\mathfrak C=(ABC)^2\), and therefore \(ABC=A\mathfrak A=B\mathfrak B=C\mathfrak C\) is necessarily an integer. Let the greatest common divisor of the numbers \(\mathfrak A,A\mathfrak A\) be \(m\), and let \(\mathfrak A=gm,\ A\mathfrak A=hm\). Then \(g\) is coprime to \(h\) and, since \(\mathfrak A\) is free of square factors, also to \(m\). Now \(hhm=gAA\mathfrak A =g\mathfrak B\mathfrak C\), so \(g\) divides \(hhm\), which is clearly impossible unless \(g=\pm1\). Thus \(\mathfrak A=\pm m,\ A=\pm h\), so it is an integer; similarly, \(B,C\) are integers. \textit{This proves the first assertion.} — Since \(\mathfrak A=BC\) has no square factors, \(B,C\) must be coprime. Similarly, \(A\) is coprime to \(C\) and to \(B\). \textit{This proves the second assertion.} — Finally, if \((\Omega)\) is satisfied by \(X=P,Y=Q\), \(Z=R\), then clearly \((\omega)\) is solved by \(x=\alpha P,y=\beta Q,z=\gamma R\). Conversely, if the latter is satisfied by \(x=p,y=q,z=r\), the former is satisfied by \(X=\beta\gamma p\), \(Y=\alpha\gamma q,Z=\alpha\beta r\). So either both are solvable or neither is. \textit{This proves the third assertion.}

\Needspace{9\baselineskip}\begin{center}Representation of zero by arbitrary ternary forms.\end{center}


\Needspace{7\baselineskip}\Article{299}
\textsc{Problem.} \textit{Given a ternary form}

\[f=axx+a'x'x'+a''x''x''+2bx'x''+2b'xx''+2b''xx'\]

\SourceLeaf{364}
\textit{determine whether it can represent zero, with values of the indeterminates that are not all \(=0\).}

\textit{Solution.} I.\ When \(a=0\), the values of \(x^\prime,x^{\prime\prime}\) can be chosen at will. The equation
\[a^\prime x^\prime x^\prime+2b x^\prime x^{\prime\prime}+a^{\prime\prime}x^{\prime\prime}x^{\prime\prime}=-2x(b^\prime x^{\prime\prime}+b^{\prime\prime}x^\prime)\]
clearly gives \(x\) a definite rational value. Whenever this gives a fraction for \(x\), we need only multiply the values of \(x,x^\prime,x^{\prime\prime}\) by its denominator to obtain integers. The only values of \(x^\prime,x^{\prime\prime}\) to exclude are those making \(b^\prime x^{\prime\prime}+b^{\prime\prime}x^\prime=0\), unless they also make \(a^\prime x^\prime x^\prime+2b x^\prime x^{\prime\prime}+a^{\prime\prime}x^{\prime\prime}x^{\prime\prime}=0\), in which case \(x\) can be chosen at will. It is also clear that all possible solutions can be obtained in this way. The case where \(b^\prime\) is zero and \(b^{\prime\prime}=0\) holds does not belong here: then \(x\) does not occur in \(f\), so \(f\) is a binary form, and whether it can represent zero must be decided from the theory of such forms applied to \(f\).

II.\ When \(a=0\) does not hold, the equation \(f=0\) is equivalent to
\[(ax+b^{\prime\prime}x^\prime+b^\prime x^{\prime\prime})^2-A^{\prime\prime}x^\prime x^\prime+2B x^\prime x^{\prime\prime}-A^\prime x^{\prime\prime}x^{\prime\prime}=0\]
where we set
\[\begin{gathered}b^{\prime\prime}b^{\prime\prime}-aa^\prime=A^{\prime\prime},\\ ab-b^\prime b^{\prime\prime}=B,\\ b^\prime b^\prime-aa^{\prime\prime}=A^\prime\end{gathered}\]
If here \(A^\prime=0\) but not \(B=0\), then choosing \(ax+b^{\prime\prime}x^\prime+b^\prime x^{\prime\prime}\) and \(x^{\prime\prime}\) at will clearly determines \(x\) and \(x^\prime\) rationally. If they are not integers, a suitable multiplier will at least make them integers. For the one value of \(x^{\prime\prime}\), namely \(x^{\prime\prime}=0\), the value of \(ax+b^{\prime\prime}x^\prime+b^\prime x^{\prime\prime}\) is not arbitrary: it too must be set \(=0\). But then \(x^\prime\) can be chosen at will and will give a rational value of \(x\). — If, however, \(A^{\prime\prime}\) is zero and \(B=0\) holds, and if \(A^\prime\) is a square \(=kk\), then \(f=0\) clearly reduces to the two linear equations
\[\begin{gathered}ax+b^{\prime\prime}x^\prime+(b^\prime+k)x^{\prime\prime}=0,\\ ax+b^{\prime\prime}x^\prime+(b^\prime-k)x^{\prime\prime}=0\end{gathered}\]
of which \textit{either} one \textit{or} the other must hold. If instead, under the same assumption, \(A^\prime\) is not a square, the solution of the proposed equation clearly depends on \(x^{\prime\prime}=0\) and \(ax+b^{\prime\prime}x^\prime=0\), which must hold \textit{simultaneously}.
\SourceLeaf{365}

It hardly needs mentioning that the method in I also applies when \(a^\prime\) is zero or \(a^{\prime\prime}=0\) holds, and the method in II when \(A^\prime=0\).

III.\ When neither \(a\) is zero nor \(A^{\prime\prime}=0\) holds, the equation \(f=0\) is equivalent to
\[A^{\prime\prime}(ax+b^{\prime\prime}x^\prime+b^\prime x^{\prime\prime})^2-(A^{\prime\prime}x^\prime-Bx^{\prime\prime})^2+Da x^{\prime\prime}x^{\prime\prime}=0\]
where \(D\) denotes the determinant of the form \(f\), or \(Da\) denotes the number \(BB-A^\prime A^{\prime\prime}\). When \(D=0\), the solution works much as at the end of the preceding case. If \(A^{\prime\prime}\) is a square \(=kk\), the proposed equation reduces to
\[\begin{gathered}kax+(kb^{\prime\prime}-A^{\prime\prime})x^\prime+(kb^\prime+B)x^{\prime\prime}=0,\\ kax+(kb^{\prime\prime}+A^{\prime\prime})x^\prime+(kb^{\prime\prime}-B)x^{\prime\prime}=0\end{gathered}\]
but if \(A^{\prime\prime}\) is not a square, we must have
\[\begin{gathered}ax+b^{\prime\prime}x^\prime+b^\prime x^{\prime\prime}=0,\\ A^{\prime\prime}x^\prime-Bx^{\prime\prime}=0\end{gathered}\]
When, however, \(D\) is not \(=0\), we have reduced the problem to
\[A^{\prime\prime}tt-uu+Davv=0\]
whose solvability can be decided by the preceding article. If its only solution is \(t=0,\ u=0,\ v=0\), then clearly the proposed equation also has no solution other than \(x=0,\ x^\prime=0,\ x^{\prime\prime}=0\). But if it has another solution, any such integer values of \(t,u,v\) give, through the equations
\[\begin{gathered}ax+b^{\prime\prime}x^\prime+b^\prime x^{\prime\prime}=t,\\ A^{\prime\prime}x^\prime-Bx^{\prime\prime}=u,\\ x^{\prime\prime}=v\end{gathered}\]
at least rational values of \(x,x^\prime,x^{\prime\prime}\). If these involve fractions, a suitable multiplier will make them integers.

As soon as \textit{one} integer solution of \(f=0\) has been found, the problem can be reduced to case I, and all solutions found as there, in the following way. Let \(f=0\) be satisfied by the values of \(x,x^\prime,x^{\prime\prime}\) given by \(\alpha,\alpha^\prime,\alpha^{\prime\prime}\), which we assume have no common factor. Choose integers \(\beta,\beta^\prime,\beta^{\prime\prime},\ \gamma,\gamma^\prime,\gamma^{\prime\prime}\) by Articles 40, 279 so that
\[\alpha(\beta^\prime\gamma^{\prime\prime}-\beta^{\prime\prime}\gamma^\prime)+\alpha^\prime(\beta^{\prime\prime}\gamma-\beta\gamma^{\prime\prime})+\alpha^{\prime\prime}(\beta\gamma^\prime-\beta^\prime\gamma)=1\]

\SourceLeaf{366}
and let \(f\) change under the substitution
\[\begin{gathered}x=\alpha y+\beta y^\prime+\gamma y^{\prime\prime}.\\ x^\prime=\alpha^\prime y+\beta^\prime y^\prime+\gamma^\prime y^{\prime\prime},\\ x^{\prime\prime}=\alpha^{\prime\prime}y+\beta^{\prime\prime}y^\prime+\gamma^{\prime\prime}y^{\prime\prime}\ldots(S)\end{gathered}\]
into
\[g=cyy+c^\prime y^\prime y^\prime+c^{\prime\prime}y^{\prime\prime}y^{\prime\prime}+2dy^\prime y^{\prime\prime}+2d^\prime yy^{\prime\prime}+2d^{\prime\prime}yy^\prime\]
Then clearly \(c=0\), and \(g\) is equivalent to \(f\). So from all solutions of \(g=0\) we readily obtain, through \(S\), all integer solutions of \(f=0\). Now I shows that all solutions of \(g=0\) are contained in the formulas
\[\begin{gathered}y=-z(c^\prime pp+2dpq+c^{\prime\prime}qq),\\ y^\prime=2z(d^{\prime\prime}pp+d^\prime pq),\\ y^{\prime\prime}=2z(d^{\prime\prime}pq+d^\prime qq)\end{gathered}\]
where \(p,q\) are arbitrary integers and \(z\) is an arbitrary number, for which fractions are also allowed, provided \(y,y^\prime,y^{\prime\prime}\) remain integers. Substituting these values of \(y,y^\prime,y^{\prime\prime}\) in \((S)\) gives all integer solutions of \(f=0\). — For example, if
\[f=xx+x^\prime x^\prime+x^{\prime\prime}x^{\prime\prime}-4x^\prime x^{\prime\prime}+2xx^{\prime\prime}+8xx^\prime\]
and one solution of \(f=0\) is \(x=1,\ x^\prime=-2,\ x^{\prime\prime}=1\), then setting \(\beta,\beta^\prime,\beta^{\prime\prime},\ \gamma,\gamma^\prime,\gamma^{\prime\prime}=0,1,0,0,0,1\) gives
\[g=y^\prime y^\prime+y^{\prime\prime}y^{\prime\prime}-4y^\prime y^{\prime\prime}+12yy^{\prime\prime}\]
So all integer solutions of \(g=0\) are contained in the formula
\[\begin{gathered}y=-z(pp-4pq+qq),\\ y^\prime=12zpq,\\ y^{\prime\prime}=12zqq\end{gathered}\]
and therefore all solutions of \(f=0\) in
\[\begin{gathered}x=-z(pp-4pq+qq)\\x^\prime=2z(pp+2pq+qq)\\x^{\prime\prime}=-z(pp-4pq-11qq)\end{gathered}\]


\Needspace{9\baselineskip}\begin{center}General solution in rational numbers of indeterminate equations of the second degree in two unknowns.\end{center}


\Needspace{7\baselineskip}\Article{300}
The problem of the preceding article directly gives the solution of the indeterminate equation
\[axx+2bxy+cyy+2dx+2ey+f=0\]

\SourceLeaf{367}
when only rational values are required. We have already solved it above, in Article 216 and the following articles, when integers are required. For all rational values of \(x,y\) can be written as \(\frac{t}{v},\frac{u}{v}\), with \(t,u,v\) integers. So solving that equation in rational numbers is clearly the same as solving
\[att+2btu+cuu+2dtv+2euv+fvv=0\]
in integers; this is the equation treated in the preceding article. Only solutions with \(v=0\) must be excluded. Such solutions cannot occur when \(bb-ac\) is not a square. For example, all rational solutions of the equation
\[xx+8xy+yy+2x-4y+1=0\]
which was solved generally in integers in Article 221, are contained in the formula
\[\begin{gathered}x=\frac{pp-4pq+qq}{pp-4pq-11qq},\\ y=-\frac{2pp+4pq+2qq}{pp-4pq-11qq}\end{gathered}\]
where \(p,q\) are any integers. — We have dealt only briefly here with these two closely connected problems, and have left out many observations related to them, partly to avoid too long a discussion and partly because we have another solution of the problem in the preceding article, based on more general principles. Since explaining it requires a deeper investigation of ternary forms, we must reserve it for another occasion.
